\section{Conclusion}

We have proposed a new method, the smeared operator product expansion (sOPE), 
to extract matrix elements from numerical nonperturbative calculations without 
power divergent mixing. The smeared operator product expansion is a general 
framework relevant to any asymptotically free theory with nonperturbative 
matrix elements that suffer from power divergent mixing. Within QCD, the most 
obvious application is to deep inelastic scattering, but other applications 
include nonperturbative determinations of $K\rightarrow \pi\pi$ decays 
\cite{Dawson:1997ic,Detmold:2005gg} and $B$-meson mixing 
\cite{Monahan:2014xra}. Beyond QCD, applications include nonperturbative 
studies 
of critical phenomena in the Heisenberg model, spin systems and other condensed 
matter systems.

In the sOPE, we expand nonlocal operators in a basis of smeared operators, 
the matrix elements of which can be determined on the lattice. We implement the 
smearing via the gradient flow, a 
classical evolution of the theory in a new dimension that smooths ultraviolet 
fluctuations. The continuum limit of these matrix elements is free of power 
divergent mixing, provided the localization scale, the smearing length, is kept 
fixed in the continuum limit. The resulting matrix elements are functions of 
two 
scales, the renormalization scale and the smearing length. The sOPE 
systematically relates these matrix elements to smeared Wilson coefficients, 
which can be calculated in perturbation theory, thereby providing a complete 
determination of the nonlocal operators.
